%---------------------------------------------------------------------
\section{Weekly Training Integration}
\label{sec:weekly}

%---------------------------------------------------------------------
\subsection{Summary of Cited Peer-Reviewed Studies}
\label{subsec:weekly_table}

\noindent\textbf{\evidencebacked} studies providing direct quantitative data for this section are listed below. Studies not independently confirmed by numeric effect-size verification are excluded from the table and flagged where used in text.

{\footnotesize
\setlength{\tabcolsep}{4pt}
\renewcommand{\arraystretch}{1.2}
\begin{longtable}{>{\raggedright\arraybackslash}p{2.2cm}
                  >{\raggedright\arraybackslash}p{2.0cm}
                  >{\raggedright\arraybackslash}p{2.6cm}
                  >{\raggedright\arraybackslash}p{2.5cm}
                  >{\raggedright\arraybackslash}p{2.8cm}
                  >{\raggedright\arraybackslash}p{1.4cm}
                  >{\raggedright\arraybackslash}p{1.2cm}}
\caption{Peer-reviewed studies informing weekly training integration}
\label{tab:weekly_studies}\\
\toprule
\textbf{Authors} & \textbf{Design / $n$} & \textbf{Protocol} & \textbf{Primary outcome} & \textbf{Effect size} & \textbf{Endpoint} & \textbf{Quality} \\
\midrule
\endfirsthead
\multicolumn{7}{l}{\small\itshape Table \ref{tab:weekly_studies} continued\ldots} \\
\toprule
\textbf{Authors} & \textbf{Design / $n$} & \textbf{Protocol} & \textbf{Primary outcome} & \textbf{Effect size} & \textbf{Endpoint} & \textbf{Quality} \\
\midrule
\endhead
\bottomrule
\endlastfoot

Stien et al.\ \cite{SO23}
& SR + meta-analysis; $n = 11$ studies (5 pooled RCTs)
& Climbing resistance training 2–3$\times$/wk, 4–10 wk vs.\ climbing-only control
& Finger strength, RFD, forearm endurance (dynamometric)
& Overall SDM = 0.57 (95\% CI 0.24–0.91); strength 0.41 (0.03–0.80); RFD 0.90 (0.21–1.61); endurance 1.23 (0.69–1.77)
& Surrogate; no on-wall gain
& GRADE: low–moderate \\

López-Rivera \& González-Badillo \cite{LRGB19}
& RCT, 3 groups; $n = 26$ advanced ($\sim$7c+/8a)
& 8 wk, 2$\times$/wk: MaxHangs ($\geq$80\% MFS) vs.\ IntHangs vs.\ Combination; all climbed 6 d/wk
& Grip endurance: time-to-failure dead hang on 11 mm edge
& IntHangs $+$45\% (ES 1.0, $p=0.002$); MaxHangs $+$34\% (ES 0.6, $p=0.010$); Combo $+$7\% (ns)
& Surrogate (dynamometric)
& GRADE: low \\

Devise et al.\ \cite{DO22}
& RCT, 4 groups; $n = 54$ experienced (41M/13F)
& 4 wk, 2$\times$/wk: F100 / F80 / F60 / control on 12 mm hold
& MFS (N), stamina, endurance
& F100 MFS $+$14.3\% ($p<0.001$); F80 MFS $+$12.4\%; F80 stamina 29.5$\to$60.4\%; F80 endurance 63.2$\to$79.2\%
& Surrogate (dynamometric)
& GRADE: low \\

Stien et al.\ \cite{SPV+21}
& RCT; $n = 16$ adv–elite (groups $n=5,6,5$)
& 5 wk, campus board: 4$\times$/wk vs.\ 2$\times$/wk vs.\ climbing-only; $\geq$48 h between sessions
& Bouldering performance, max force, campus moves, RFD
& 4$\times$/wk: force/RFD $\uparrow$ vs.\ control; 2$\times$/wk: bouldering/campus $\uparrow$ vs.\ control; no between-freq.\ difference
& Mixed (bouldering: functional; force/RFD: surrogate)
& GRADE: very low \\

Wilson et al.\ \cite{WMR+12}
& Meta-analysis; 21 RCTs/CCTs; $n = 422$ ES
& Concurrent strength + endurance vs.\ strength-only; running/cycling; varied freq.\ and duration
& Hypertrophy, maximal strength, power
& Strength ES: 1.76 (solo) vs.\ 1.44 (concurrent, $-$18\%); hypertrophy $-$31\%; power $-$40\%; endurance duration $r=-0.29$ to $-0.75$
& Surrogate; non-climbing; inferential transfer
& GRADE: moderate \\

\end{longtable}
}% end footnotesize

%---------------------------------------------------------------------
\subsection{Fatigue Interference Between Modalities}
\label{subsec:fatigue_interference}

\begin{evidencebox}{Supplemental Finger Training Outperforms Climbing Alone (Stien 2023 meta-analysis) · \evidencebacked · ESS 7}
Supplemental finger-strength training (hangboard or campus board) superimposed on regular climbing training produces significantly greater gains in climbing-specific tests than climbing alone: overall meta-analytic SDM = 0.57 (95\% CI 0.24–0.91), finger strength SDM = 0.41 (0.03–0.80), forearm endurance SDM = 1.23 (0.69–1.77) (Stien et al.\ 2023, \textit{Biol Sport} 40(1):179–191; \cite{SO23}). \textit{Endpoint: all pooled outcomes are dynamometric surrogates (force, endurance, RFD); no included trial showed significant improvement on an on-wall climbing test.} These studies implicitly impose $\geq$48 h between sessions by scheduling 2–3 sessions/wk but do not experimentally vary session order or inter-session gap. No published RCT has directly randomized the temporal proximity between hangboard training and limit bouldering to measure interference on adaptation outcomes in climbers.

The concurrent-training interference effect (Wilson et al.\ 2012) is well-characterised for lower-body strength plus cardiovascular endurance: adding endurance training at high frequency ($r = -0.35$) or duration ($r = -0.75$) significantly attenuates strength, hypertrophy, and power gains versus strength-only training, with running producing larger deficits than cycling. The molecular basis is AMPK-mediated inhibition of mTORC1: AMPK requires $\geq$3 h post-endurance to return to baseline, while mTORC1 remains elevated $\approx$18 h post-resistance loading.
\end{evidencebox}

\begin{bioplausbox}{Concurrent Training Interference: Peripheral and Molecular Mechanisms · \bioplausible · ESS 2}
Limit bouldering ($\geq$90–100\% effort) and hangboard max-hang sessions share the same primary effector tissue: the finger flexor musculotendinous unit (FDP/FDS) and annular pulley system (A2, A4). Both impose near-maximal isometric loading on a 6–14 mm crimp, activating the same motor-unit pool at high discharge rates. The dominant interference mechanism is not the classical AMPK/mTOR competition (which is driven by systemic aerobic endurance), but rather: (1) acute peripheral fatigue — H\textsuperscript{+} accumulation and impaired sarcoplasmic Ca\textsuperscript{2+} release reduce peak force capacity; (2) phosphocreatine depletion — PCr returns to $\approx$50\% of resting levels within 2–4 min but full resynthesis takes $\approx$8–12 min post-maximal isometric effort; (3) central inhibition — reduced corticospinal drive via Golgi tendon organ (GTO) afference after maximal isometric loading. A second max-effort session initiated before full neuromuscular recovery delivers a degraded force stimulus, reducing the adaptive signal for maximal strength.

The AMPK/mTOR model applies most cleanly to combining route climbing endurance (aerobic forearm demand) with hangboard or limit bouldering (maximal isometric). Climbing endurance is a local aerobic-to-anaerobic continuum without the eccentric locomotion that drives the strongest interference in Wilson et al.\ (2012), so the magnitude of molecular interference in climbers is likely $>0$ but smaller than the $-$18\% strength deficit observed in running-concurrent protocols.

\textit{Evidence gap: no RCT has measured post-hangboard neuromuscular recovery time in the finger flexors at a temporal resolution sufficient to guide programming, and no study has measured AMPK or mTOR pathway activity in finger flexors following climbing-specific endurance versus limit bouldering.}
\end{bioplausbox}

\begin{folklorebox}{48\,h Minimum Between High-Intensity Finger Sessions · \folkloreverified · ESS 1}
\textbf{Source: Lattice Training blog (Randall \& Torr); Anderson \& Anderson, \textit{The Rock Climber's Training Manual}, 2nd ed., 2014; Hörst, \textit{Training for Climbing}, 3rd ed., 2016; r/climbharder wiki.}

Practitioner consensus across all major coaching frameworks prescribes $\geq$48 h between any two sessions imposing high-intensity finger loading (max-hangs, campus board, limit bouldering, Moonboard benchmarks). Some coaches (Neil Gresham, UKClimbing; Power Company Podcast) recommend 72 h between max-intensity finger sessions for advanced climbers. Both the 48 h and 72 h thresholds are study-design conventions embedded in climbing RCTs as methodological safeguards (Stien 2021 specified $\geq$48 h between campus sessions), not empirically optimised prescriptions.
\end{folklorebox}

\subsubsection*{Moonboard Sessions}

\begin{bioplausbox}{Moonboard as High-Intensity Finger Session: Scheduling Equivalence · \bioplausible · ESS 2}
The Moonboard is a 40$^\circ$-overhung standardised board using benchmark problems with $\leq$1-phalange-depth holds. A full session (60–90 min, $\geq$8–10 max-effort attempts per problem) involves near-maximal finger loading, high CNS demand, and significant metabolic depletion — a fatigue profile overlapping limit bouldering on a steep wall, not top-rope endurance. Moonboard sessions should be categorised as high-intensity finger sessions equivalent to limit bouldering for scheduling purposes; the $\geq$48 h separation constraint applies identically. Placing a Moonboard session adjacent to a hangboard max-strength day (e.g., Moonboard Monday, Hangboard Tuesday) risks the same attenuation of MFS stimulus quality as doing limit bouldering the day before max-hangs.

\textit{Evidence gap: no study has characterised post-Moonboard recovery of MFS, RFD, or tendon stiffness, nor compared its fatigue profile to an equivalent-duration limit-bouldering session.}\singlesource
\end{bioplausbox}

\subsubsection*{Top-Rope and Lead Endurance Sessions}

\begin{bioplausbox}{Route Endurance Fatigue Profile: Lower Interference Than Limit Bouldering · \bioplausible · ESS 2}
Top-rope and lead climbing at 60–75\% maximum effort impose forearm endurance demand (aerobic-glycolytic) with substantially lower peak finger force per hold than limit bouldering (estimated 50–70\% MFS vs.\ $\geq$90\% MFS). Active recovery research in climbers demonstrates that low-to-moderate climbing as an inter-session recovery modality accumulates forearm fatigue relative to non-climbing locomotion (performance impulse 29.5\% lower after climbing vs.\ 18.4\% lower after walking, $p = 0.009$, $d = 0.54$).\singlesource The fatigue profile of route endurance (sustained aerobic demand at submaximal grip force) is qualitatively distinct from max-hang fatigue and is unlikely to degrade max-strength stimulus quality as severely as placing a Moonboard session adjacent to a hangboard day. Optimal placement: route endurance sessions should bookend the high-intensity block, either as an early-week warm-up modality or as the final training day before a rest day. They should \emph{not} immediately precede a hangboard max-strength session.
\end{bioplausbox}

%---------------------------------------------------------------------
\subsection{Session Sequencing Rules}
\label{subsec:sequencing}

\begin{bioplausbox}{Hangboard Before Bouldering for Max-Strength Priority · \bioplausible · ESS 2}
\textbf{Hangboard before or after limit bouldering?}

Devise et al.\ (2022) demonstrate intensity-specific adaptation: F100 (100\% MFS) significantly improves MFS but not stamina or endurance; F80 (80\% MFS) improves MFS, stamina, and endurance; F60 (60\% MFS) improves stamina and endurance but not MFS. This specificity has direct sequencing implications: max-hang adaptation requires delivery of the stimulus at close to peak force output. Any prior modality that depletes finger-flexor force capacity shifts the effective stimulus downward — an intended F100 session becomes an F80 or F60 session after a 2–3 h limit bouldering session (estimated MVC decrement 10–25\%, inferred from general post-climbing strength decrements; no direct hangboard-after-bouldering measurement exists in the peer-reviewed literature).

\textbf{Principle:} When the goal is maximal finger strength or RFD, hangboard max-hangs must precede — not follow — limit bouldering, and should be executed at session start after warm-up. When the goal is endurance (repeaters at F60–F80 intensity), this constraint is weaker because the submaximal stimulus is less sensitive to pre-fatigue.

\textit{Evidence gap: no RCT has randomised hangboard-before vs.\ hangboard-after bouldering to directly test the sequence effect in climbers.}
\end{bioplausbox}

\begin{bioplausbox}{Same-Day Hangboard + Limit Bouldering: Feasibility Conditions · \bioplausible · ESS 2}
\textbf{Same-day hangboard + limit bouldering.}

Same-session combination is feasible if: (a) max-strength hangboard is performed first, immediately after warm-up; (b) within-session passive gap of approximately 10--25\,min (broad range; no climbing data; general PCr resynthesis literature suggests $\approx$8--12\,min for full restoration post-maximal isometric effort) provides partial PCr recovery without additional forearm loading; (c) total hangboard volume reduced approximately 30--60\% versus a dedicated session (a pragmatic range reflecting the absence of dose-response climbing data; individual tolerance varies); (d) total high-intensity time-under-tension ideally limited to roughly 45--90\,s across both modalities to moderate AMPK activation (these values are ESS~2 extrapolations from concurrent-training literature, not empirically derived thresholds for climbing). Combining both modalities in one session does not substitute for two separated sessions, as cumulative volume limits the training signal of each.\singlesource

Sequential strength-before-endurance within a concurrent session preserves the quality of the strength stimulus in general concurrent-training literature; the same principle applies by bioplausible transfer to upper-limb isometric loading.
\end{bioplausbox}

\begin{folklorebox}{Hangboard-First Session Ordering · \folkloreverified · ESS 1}
\textbf{Source: Eva López periodization blog; Lattice Training blog; Anderson \& Anderson 2014.}

Eva López's periodization guidelines explicitly recommend hangboard work at the \emph{start} of a session when freshness is required for max-effort protocol quality. The Anderson brothers and Lattice Training both advocate same-day hangboard-first ordering when schedule constraints prevent separate days. These are practitioner heuristics without RCT validation.
\end{folklorebox}

%---------------------------------------------------------------------
\subsection{Example Weekly Structures}
\label{subsec:templates}

\noindent\textbf{Scope:} All templates target intermediate-to-advanced climbers (V6–V9 / 5.12–5.13c redpoint), 4–5 training days/wk. Intensity descriptors: \textit{max} = $\geq$90–100\% MFS; \textit{high submaximal} = 75–89\% MFS; \textit{submaximal} = 60–74\% MFS; \textit{technique/recovery} = $\leq$60\%, skill-focused. No complete 7-day weekly structure has been tested as a unit in a climbing-specific RCT; component protocols carry their individual evidence ratings.

\subsubsection*{Template A — Strength-Priority Block (4 days/wk, hangboard-led)}

\noindent\textit{Primary goal: maximal finger strength and RFD. Most directly evidence-proximate template; hangboard protocol matches the 2$\times$/wk frequency used in López-Rivera \& González-Badillo (2019) and consistent with Devise et al.\ (2022) F100/F80 parameters.}

\begin{itemize}[leftmargin=1.6cm,labelwidth=1.2cm,align=left]
  \item[\textbf{Mon:}] \textbf{Hangboard max-hangs.} $\geq$80\% MFS, 10 s hangs, 3–5 sets $\times$ 3–5 reps, 3–5 min inter-set rest, 18–20 mm edge; total $\approx$20–30 min. \textbf{Evidence-backed:} López-Rivera \& González-Badillo (2019) — MaxHangs protocol structure and intensity; Devise et al.\ (2022) — F100/F80 intensity parameters.
  \item[\textbf{Tue:}] \textbf{Rest or non-climbing antagonist work} (wrist extensors, shoulder external rotation), $\leq$30 min. \textbf{Bioplausible:} passive recovery to allow neuromuscular restoration; 24 h post-hangboard likely insufficient for full MVC recovery.
  \item[\textbf{Wed:}] \textbf{Limit bouldering / Moonboard projecting.} 90–100\% effort; 2–4 attempts per project problem; 5–10 min rest between attempts; climbing time $\approx$60–90 min; total session $\approx$2–2.5 h including warm-up. \textbf{Folklore:} Anderson \& Anderson 2014, r/climbharder wiki — structure of projecting sessions; no RCT specifies volume/rest parameters for limit bouldering as a training modality.
  \item[\textbf{Thu:}] \textbf{Rest.} $\geq$48 h separation from Wednesday's high-intensity finger session before next max-load day. \textbf{Bioplausible.}
  \item[\textbf{Fri:}] \textbf{Hangboard max-hangs} (second session; same as Monday, or Moonboard benchmark problems at 75–85\% effort, $\approx$90 min if goal shifts to power-endurance). \textbf{Evidence-backed} for hangboard frequency (2$\times$/wk) per Stien et al.\ (2021) and López-Rivera \& González-Badillo (2019). Moonboard substitution: \textbf{Folklore} (Lattice Training periodization framework).
  \item[\textbf{Sat:}] \textbf{Top-rope/lead endurance.} Routes at 60–75\% max effort; $\approx$2 h; $\geq$10 routes or equivalent mileage. \textbf{Bioplausible:} aerobic forearm demand stimulates capillarisation and lactate-clearance adaptations; low enough intensity not to significantly attenuate next-week hangboard quality.
  \item[\textbf{Sun:}] \textbf{Rest.}
\end{itemize}

\noindent\textit{Classification: Hangboard days are the most directly evidence-backed component. The overall 7-day arrangement is folklore-derived; the sequencing logic (hard finger $\to$ rest $\to$ limit bouldering $\to$ rest $\to$ hangboard $\to$ endurance $\to$ rest) is bioplausible from concurrent-training principles.}

\subsubsection*{Template B — Bouldering-Led Block (4 days/wk, board-climbing priority)}

\noindent\textit{Primary goal: project-specific bouldering performance and on-wall power. Hangboard is maintained at reduced volume as a supplemental stimulus.}

\begin{itemize}[leftmargin=1.6cm,labelwidth=1.2cm,align=left]
  \item[\textbf{Mon:}] \textbf{Limit bouldering.} 90–100\% effort; project-style attempts; 2–2.5 h total (30 min warm-up, 60–90 min projecting). \textbf{Folklore:} Anderson \& Anderson 2014 ``Limit Bouldering'' chapter; r/climbharder wiki.
  \item[\textbf{Tue:}] \textbf{Hangboard max-hangs (reduced volume).} $\geq$80\% MFS; 10 s, 3 $\times$ 4–5 reps; performed 24–48 h post-bouldering — a scheduling compromise; volume reduced 20–30\% vs.\ a fully recovered session if within 48 h. \textbf{Bioplausible:} sub-optimal recovery window; \textbf{Folklore:} Lattice Training blog explicitly flags the 24 h post-bouldering hangboard trade-off.
  \item[\textbf{Wed:}] \textbf{Rest or antagonist/mobility.} \textbf{Bioplausible.}
  \item[\textbf{Thu:}] \textbf{Moonboard benchmark problems.} 70–85\% effort; emphasis on movement quality over max-effort; $\approx$90 min. \textbf{Folklore:} Moonboard app programming, r/climbharder.
  \item[\textbf{Fri:}] \textbf{Rest.} $\geq$48 h before next high-intensity day. \textbf{Bioplausible.}
  \item[\textbf{Sat:}] \textbf{Top-rope/lead endurance or ARC.} 60–75\% max effort, $\approx$2 h. \textbf{Folklore:} Anderson \& Anderson 2014 — ARC training definition and prescription.
  \item[\textbf{Sun:}] \textbf{Rest.}
\end{itemize}

\noindent\textit{Classification: Folklore-derived overall structure; individual hangboard parameters are evidence-backed. The Tuesday hangboard after Monday bouldering is a pragmatic scheduling compromise, classified as bioplausible rather than optimal.}

\subsubsection*{Template C — Evidence-Constrained Balanced Block (4 days/wk)}

\noindent\textit{Primary goal: balanced finger-strength and endurance development. This template most closely replicates the 2$\times$/wk hangboard-plus-normal-climbing design used across Hermans et al.\ (2022)\footnote{Hermans et al.\ (2022) \textit{Front Sports Act Living} 4:888158. \cite{HO22}. RCT; $n = 35$; 10 wk, 2 sessions/wk (48 min/session), Beastmaker 1000 app-guided protocol; $F_\mathrm{peak}$ $\uparrow$ ($p < 0.05$, large within-group ES) vs.\ control. GRADE: low. \toverify: cited by fewer than 3 independent sources; requires verification. See Appendix~\ref{app:toverify}, entry~2.} and consistent with López-Rivera \& González-Badillo (2019).}

\begin{itemize}[leftmargin=1.6cm,labelwidth=1.2cm,align=left]
  \item[\textbf{Mon:}] \textbf{Hangboard protocol} (2$\times$/wk schedule; Beastmaker-style or equivalent force-quantified protocol; 48 min including warm-up; progressive overload weekly). \textbf{Evidence-backed:} 2$\times$/wk frequency consistent with López-Rivera \& González-Badillo (2019) and Devise et al.\ (2022) intervention structures.
  \item[\textbf{Tue:}] \textbf{Rest.} $\geq$48 h minimum between hangboard sessions. \textbf{Evidence-backed:} study constraint in climbing-specific RCTs (Stien 2021; Devise 2022).
  \item[\textbf{Wed:}] \textbf{Limit bouldering or Moonboard.} 90–100\% effort; 60–90 min climbing; equivalent to ``normal climbing'' continued by participants in the intervention arms of the cited RCTs. \textbf{Folklore} for specific volume/intensity; consistent with the design intent of the cited trials.
  \item[\textbf{Thu:}] \textbf{Rest or antagonist work.} \textbf{Bioplausible.}
  \item[\textbf{Fri:}] \textbf{Hangboard protocol} (second session of week; identical structure to Monday). \textbf{Evidence-backed:} 2$\times$/wk schedule per intervention design.
  \item[\textbf{Sat:}] \textbf{Route climbing/lead endurance.} 60–75\% max effort; $\approx$2 h. \textbf{Bioplausible:} aerobic demand modality separate from finger-strength focus days.
  \item[\textbf{Sun:}] \textbf{Rest.}
\end{itemize}

\noindent\textit{Classification: The hangboard days are the most directly evidence-backed element of any template. The specific 7-day composition is folklore-derived, but internally consistent with the evidence-based principle of $\geq$48 h between high-intensity finger sessions.}

\subsubsection*{General Constraint Across All Templates}

\begin{folklorebox}{Weekly Hard Finger Exposure Cap: 2--3 Sessions · \folkloreverified · ESS 1}
\textbf{Source: Lattice Training; Anderson \& Anderson 2014; r/climbharder wiki.}

Keep total ``hard finger exposures'' (limit bouldering, Moonboard, max-hangs at $\geq$80\% MFS) at \textbf{2–3 per week} for most V6–V9 climbers unless individual recovery capacity has been demonstrated to support higher frequency. This cap is a widespread coaching heuristic with no RCT dose-response validation in the target population.
\end{folklorebox}

%---------------------------------------------------------------------
\subsection{Deload Triggers and Protocol}
\label{subsec:deload}

\begin{bioplausbox}{Overreaching Taxonomy and Deload Indicators for Climbers · \bioplausible · ESS 2}
\textbf{Overreaching taxonomy (from non-climbing sports literature; bioplausible transfer).}

Functional overreaching (FOR) = planned short-term performance decrement with supercompensation on recovery, full restoration $\leq$2 wk. Non-functional overreaching (NFOR) = persistent decrements requiring $\leq$4 wk to resolve. Overtraining syndrome (OTS) = months-to-years recovery. FOR is tolerated training stress; NFOR is the injury-adjacent state. Resting HR elevation of 3–5 bpm above baseline with HRV recovering within 24–72 h after a rest day indicates FOR; persistent HR elevation $>$5–10 bpm above baseline on rest days with declining HRV over $\geq$5 consecutive training days indicates NFOR.

Deload indicators for climbing application (ESS~2 — all values extrapolated from general resistance-training literature; none validated in climbing-specific studies): (1) max-hang force output declining roughly $>$5--15\% over 2--3 consecutive sessions vs.\ 4-week rolling average (a plausible signal range; the exact threshold of ``10\%'' carries false precision and should be treated as an approximate sentinel); (2) subjective RPE for a fixed hangboard protocol increasing $\geq$1--2 points (0--10 scale) without load increase (this range reflects inter-individual RPE variability; a single 1.5-point increase is insufficient signal in isolation); (3) finger/pulley pain on VAS $\geq$3/10 at rest or $\geq$5/10 during loading — \textbf{mandatory hard stop} regardless of other indicators (may signal A2/A4 pulley irritation requiring medical evaluation; this threshold is the one most directly supported by injury-risk rationale, even though no RCT has validated the specific VAS cut-point); (4) sustained HRV (RMSSD) downward trend over $\geq$4--7 consecutive training days combined with any two of markers (1), (2), or (5) (the ``5 days'' figure is a convention from endurance coaching, not a validated climbing threshold); (5) warm-up grades feeling $\geq$1 full grade harder than baseline for $\geq$5--10 days.

Connective tissue (pulley, tendon) adapts more slowly than muscle ($\approx$12 wk for collagen turnover); neuromuscular indicators may lag soft-tissue overload, making subjective fatigue an imperfect sentinel for pulley injury risk.

\textit{Evidence gap: no climbing-specific study has validated any HRV threshold, RMSSD trajectory, or RPE criterion for deload timing in climbers; the thresholds above are extrapolated from general resistance-training literature.}
\end{bioplausbox}

\begin{bioplausbox}{Deload Protocol: 40--60\,\% Volume Reduction for 7 Days · \bioplausible · ESS 2}
\textbf{Deload protocol.} Reduce total volume 40–60\% across all modalities simultaneously for 7 days. Maintain $\geq$1 high-intensity exposure (e.g., 2–3 sets of max-hangs at $\geq$80\% MFS) to preserve neural qualities. Eliminate repeated-failure pump sessions and any session performed to forearm failure. Simultaneous reduction of all modalities restores AMPK/mTOR balance faster than isolated forearm rest when route climbing has been part of the loading block.
\end{bioplausbox}

\begin{folklorebox}{Planned Deload Every 3--6 Weeks · \folkloreverified · ESS 1}
\textbf{Source: Lattice Training (Progression Pyramid framework); Anderson \& Anderson 2014; Hörst 2016; r/climbharder wiki.}

Practitioner consensus recommends a planned deload (50–60\% volume reduction, maintained intensity $\geq$70\% MFS) every 3–6 weeks: Lattice schedules a recovery week after every 3-week loading block; the Anderson brothers specify a rest week every 4th week. Additional folklore triggers: Moonboard benchmark send rate halving for 2 consecutive weeks at the same effort level; inability to complete $>$80\% of planned hard moves for 2 consecutive sessions. Neither schedule nor trigger threshold has been tested in a climbing-specific RCT.
\end{folklorebox}

%---------------------------------------------------------------------
\subsection{What the Literature Cannot Yet Answer}
\label{subsec:open_questions}

The following gaps constitute genuine absences in the peer-reviewed climbing literature as of April 2026.

\begin{enumerate}

  \item \textbf{Optimal inter-session recovery interval.} No RCT has compared 24 h vs.\ 48 h vs.\ 72 h spacing between two high-intensity finger sessions ($\geq$80\% MFS hangboard or limit bouldering) in trained climbers, measuring MFS, RFD, and tendon morphology (ultrasound cross-section) as endpoints. The 48 h heuristic embedded in study designs is a methodological convention, not an empirically optimised prescription.

  \item \textbf{Session sequencing effect on long-term adaptation.} No RCT has randomised hangboard-before vs.\ hangboard-after bouldering (same day or adjacent days) and measured 8–12 week finger-strength or climbing-performance outcomes with matched total weekly volume. All current sequencing rules are extrapolated from concurrent-training principles and practitioner experience.

  \item \textbf{Moonboard-specific fatigue profile and interference.} No study has characterised post-Moonboard recovery of MFS or tendon stiffness, quantified the interference between Moonboard sessions and hangboard max-hangs within the same microcycle, or compared the Moonboard's fatigue profile to equivalent-duration limit bouldering on a less-steep wall.

  \item \textbf{Concurrent-training interference magnitude in climbing.} No RCT has measured AMPK or mTOR pathway activity in finger flexors following route-climbing endurance versus limit bouldering, nor directly compared adaptation outcomes in climbers who separate these modalities by 6 h, 24 h, or 48 h. The application of Wilson et al.\ (2012) to isometric forearm loading remains inferential.

  \item \textbf{Maintenance frequency during redpointing phases.} Stien et al.\ (2021) compared 2$\times$/wk vs.\ 4$\times$/wk campus board for acquisition; no study has examined 1$\times$/wk vs.\ 2$\times$/wk maintenance of finger-strength adaptations during a concurrent limit-bouldering and route-climbing performance phase.

  \item \textbf{Validated readiness and deload-timing metrics.} No climbing-specific study has validated any HRV threshold, performance-decrement criterion, or multivariate readiness algorithm that reliably predicts functional overreaching or pulley-injury risk in V6–V9 climbers undergoing concurrent modality programmes. Individual pulley load tolerance likely varies substantially by training age, sex, and tissue stiffness — moderating variables that have not been systematically quantified.

\end{enumerate}
